\appsection{MCMC Estimation and Marginal Likelihood Details}
\label{app:mcmc_details}
%==============================================================================

All three models are estimated using a Metropolis-within-Gibbs sampler. The
parameter vector is partitioned into blocks, with each subject's structural
parameters forming one block and each pair of hyperparameters (population mean
and dispersion for each quantity) forming another. Each block is updated
sequentially using a Metropolis-Hastings step with an adaptive multivariate
normal proposal distribution. The proposal covariance is estimated from the
empirical covariance of an A/B chain split: Group A chains use the running
covariance of Group B to construct proposals, and vice versa. This coupled
design improves mixing while maintaining detailed balance within each group.

For each model, 2000 chains are run in parallel on GPU (1000 per A/B group).
Every 10th iteration is retained (thinning factor of 10), and the first half
of the retained samples is discarded as burn-in. Table~\ref{tab:mcmc_summary}
summarizes the sampling configuration and convergence diagnostics for each
model.

\begin{table}[h!]
  \centering
  \begin{tabular}{lccc}
    \toprule
    & Model 1 & Model 2 & Model 3 \\
    & (CRRA + Prelec) & (CRRA + GE) & (CRRA + Power) \\
    \midrule
    Structural parameters per subject & 4 & 4 & 3 \\
    Predictive quantities & 4 & 4 & 4 \\
    Hyperparameters & 8 & 8 & 8 \\
    Total parameters & 452 & 452 & 341 \\
    Chains (A + B) & 2000 & 2000 & 2000 \\
    Collection iterations & 6000 & 6000 & 10000 \\
    Thinning & 10 & 10 & 10 \\
    Retained samples per chain & 600 & 600 & 1000 \\
    Post burn-in samples per chain & 300 & 300 & 500 \\
    Total posterior draws & 600{,}000 & 600{,}000 & 1{,}000{,}000 \\
    Median $\hat{R}$ & \py{diag_rhat_med['crra']} & \py{diag_rhat_med['ge']} & \py{diag_rhat_med['power']} \\
    Max $\hat{R}$ & \py{diag_rhat_max['crra']} & \py{diag_rhat_max['ge']} & \py{diag_rhat_max['power']} \\
    Mean acceptance rate & \py{diag_accept['crra']}\% & \py{diag_accept['ge']}\% & \py{diag_accept['power']}\% \\
    Time per iteration (s) & \pending{crra_time_per_iter} & \pending{ge_time_per_iter} & \pending{power_time_per_iter} \\
    \bottomrule
  \end{tabular}
  \caption{MCMC sampling configuration and convergence diagnostics.}
  \label{tab:mcmc_summary}
\end{table}

All three models share the CRRA utility function, so their per-iteration cost
is similar, governed by the analytic Jacobian decomposition for Models 1 and 2
and the Gram determinant $\sqrt{\det(J^\top J)}$ for Model~3. Model~3 is the
fastest per iteration because it has the fewest structural parameters and uses
the analytic CRRA inverse.

The split $\hat{R}$ statistic is computed for every subject-parameter
combination and every hyperparameter, using the second half of the retained
samples split again into two segments. Elevated $\hat{R}$ values are
concentrated in the noise parameter $\lambda$ for subjects whose choice data
are consistent with multiple modes in the posterior. These subjects exhibit
bimodal CE posteriors, as discussed in Section \ref{sec:estimation}.

\subsection{Marginal Likelihood Computation}

The marginal likelihoods reported in Section~\ref{sec:model_comparison} are
computed with three estimators, each operating on the posterior draws
described above.

The warp bridge sampling estimator follows \textcite{MengWong1996}, with the
warp transformation of \textcite{MengSchilling2002} applied to the posterior
draws so that the proposal distribution and the warped posterior overlap. The
proposal is \pending{bridge_proposal_detail}, and the Meng-Wong fixed point
iteration is run for \pending{bridge_n_iterations} iterations. Standard
errors come from the standard asymptotic formula for the bridge estimator.
Because the bridge operates on the full parameter space, the proposal must
cover 452 dimensions for Models 1 and 2 and 341 for Model 3, which makes
overlap with the posterior the binding practical constraint.

The Rao-Blackwellized bridge addresses that constraint by integrating out the
subject-level parameters before bridging. For each hyperparameter draw, the
per-subject integrated likelihood is estimated by
\pending{rb_subject_integration_method}, reducing the bridge to the
8-dimensional hyperparameter space, where a moment-matched proposal overlaps
the marginal posterior easily.

Thermodynamic integration runs independent MCMC at each rung of a
temperature ladder of \pending{n_temperatures} temperatures with
\pending{ti_schedule} spacing, raising the likelihood to the power
$\beta_t$ while keeping the prior fixed. The headline estimate uses the
stepping-stone estimator of \textcite{Xie2011}, with the trapezoidal
estimate along the path sampling identity of \textcite{GelmanMeng1998}
serving as an internal consistency check. The per-model compute cost of the
full ladder is governed by the same per-iteration costs reported in
Table~\ref{tab:mcmc_summary}.
